\chapter{Conclusiones y posibles ampliaciones}

\section{Cumplimiento de objetivos}

El trabajo desarrollado cumple el objetivo principal: construir una fase del
laboratorio virtual capaz de transformar descripciones en lenguaje natural en
modelos de decisión ejecutables. Para ello se diseñó un pipeline con agentes
especializados, revisión humana, persistencia de artefactos, memoria de
conocimiento y validación mediante pruebas.

Las contribuciones principales son:

\begin{itemize}
\item un Router que coordina el flujo completo del pipeline;
\item agentes especializados para investigación, formalización, razonamiento y
construcción;
\item una memoria persistente basada en PostgreSQL, MinIO, Neo4j y Qdrant;
\item un contrato de modelo sencillo e integrable con la fase de simulación;
\item una metodología de desarrollo apoyada en coding agents y criterios de
aceptación.
\end{itemize}

\section{Conclusiones técnicas}

La principal conclusión técnica es que los modelos de lenguaje son más útiles
cuando se integran en un sistema con límites claros. Un agente aislado puede
generar resultados convincentes, pero un sistema de desarrollo necesita
artefactos, validación, memoria y revisión.

La segunda conclusión es que la memoria no debe ser un almacén indiscriminado.
Guardar solo artefactos aceptados permite que ejecuciones futuras se apoyen en
conocimiento curado.

La tercera conclusión es metodológica: los agentes de programación son una
herramienta muy potente para implementar, probar y refactorizar, pero las
decisiones de diseño siguen necesitando contexto humano.

\section{Posibles ampliaciones}

Entre las líneas futuras destacan:

\begin{itemize}
\item añadir más paradigmas de decisión y experimentos de validación;
\item medir cuantitativamente la calidad de los modelos generados;
\item mejorar las métricas de recuperación de memoria;
\item ampliar el sistema de detección de contradicciones;
\item integrar explicaciones más detalladas de los modelos en la fase de
análisis e informes;
\item automatizar más partes de la planificación con Linear manteniendo revisión
humana en decisiones de arquitectura.
\end{itemize}
